%% ============================================================
%%  Jurnal Ilmiah — Sistem Informasi Absensi Digital Berbasis
%%  GPS, QR Code, dan Kamera dengan Gamifikasi serta Integrasi
%%  E-Learning EduRAG pada SMAN 2 Surabaya
%%
%%  Target  : Jurnal Nasional Terakreditasi SINTA 1-2 /
%%            Internasional Scopus Q2
%%  Template: Article standar (tidak butuh elsarticle)
%%  Compile : pandoc jurnal_sman2sby.tex --bibliography=references_sman2sby.bib
%%            --citeproc -o jurnal_sman2sby.docx
%% ============================================================

\documentclass[12pt,a4paper]{article}

\usepackage[T1]{fontenc}
\usepackage[utf8]{inputenc}
\usepackage[bahasa]{babel}
\usepackage{lmodern}
\usepackage{microtype}
\usepackage[top=3cm,bottom=3cm,left=3cm,right=3cm]{geometry}
\usepackage{setspace}
\doublespacing
\usepackage{amsmath,amssymb}
\usepackage{booktabs}
\usepackage{tabularx}
\usepackage{multirow}
\usepackage{longtable}
\usepackage{array}
\usepackage{graphicx}
\usepackage{tikz}
\usetikzlibrary{shapes.geometric,arrows.meta,positioning,fit,backgrounds}
\usepackage{pgfplots}
\pgfplotsset{compat=1.18}
\usepackage{listings}
\usepackage{xcolor}
\usepackage[hidelinks]{hyperref}
\usepackage{url}
\usepackage{caption}
\usepackage{subcaption}
\usepackage{enumitem}
\usepackage{fancyhdr}
\usepackage{titlesec}
\usepackage{abstract}
\usepackage{cite}
\usepackage{natbib}

%% ── Tikz styles ───────────────────────────────────────────────
\tikzset{
  kotak/.style={rectangle,draw,fill=blue!8,rounded corners=4pt,
                text width=3cm,align=center,minimum height=0.9cm,font=\small},
  kotakw/.style={rectangle,draw,fill=blue!8,rounded corners=4pt,
                text width=6cm,align=center,minimum height=0.9cm,font=\small},
  kotakr/.style={rectangle,draw,fill=red!8,rounded corners=4pt,
                text width=3cm,align=center,minimum height=0.9cm,font=\small},
  kotakg/.style={rectangle,draw,fill=green!8,rounded corners=4pt,
                text width=3cm,align=center,minimum height=0.9cm,font=\small},
  panah/.style={-{Latex[length=2mm]},thick},
  db/.style={cylinder,draw,fill=orange!10,shape border rotate=90,
             text width=2.4cm,align=center,font=\scriptsize,minimum height=0.8cm}
}

%% ── Header/footer ─────────────────────────────────────────────
\pagestyle{fancy}
\fancyhf{}
\rhead{\small Sistem Informasi Sekolah Berbasis AI — SMAN 2 Surabaya}
\lhead{\small Jurnal Teknologi Pendidikan}
\cfoot{\thepage}

%% ── Section formatting ────────────────────────────────────────
\titleformat{\section}{\large\bfseries}{{\thesection}.}{0.5em}{}
\titleformat{\subsection}{\normalsize\bfseries}{{\thesubsection}}{0.5em}{}
\titleformat{\subsubsection}{\normalsize\itshape}{{\thesubsubsection}}{0.5em}{}

%% ── Custom commands ───────────────────────────────────────────
\newcommand{\sman}{SMAN~2~Surabaya}
\newcommand{\edurag}{EduRAG}
\newcommand{\sistem}{\textit{SiAbsensi}}
\definecolor{codebg}{RGB}{248,248,248}
\lstset{backgroundcolor=\color{codebg},basicstyle=\ttfamily\small,
        frame=single,breaklines=true,captionpos=b}

%% ============================================================
\begin{document}

%% ── Halaman Judul ─────────────────────────────────────────────
\begin{titlepage}
\centering
\vspace*{2cm}
{\LARGE\bfseries Sistem Informasi Absensi Digital Berbasis GPS,
QR~Code, dan Kamera dengan Gamifikasi serta Integrasi
Platform E-Learning \edurag{} pada \sman{}\par}
\vspace{1cm}
{\large\itshape Digital Attendance Information System Based on GPS,
QR~Code, and Camera with Gamification and Integration of
\edurag{} E-Learning Platform at SMAN~2~Surabaya\par}
\vspace{2cm}
{\large
Indri Rahmawati\textsuperscript{1*},
Budi Santoso\textsuperscript{2},
Sari Wijayanti\textsuperscript{1},
Andi Nugroho\textsuperscript{3}
\par}
\vspace{0.5cm}
{\normalsize
\textsuperscript{1}Program Studi Teknik Informatika, Universitas
Airlangga, Surabaya 60115, Indonesia\\
\textsuperscript{2}Jurusan Ilmu Komputer, Institut Teknologi
Sepuluh Nopember, Surabaya 60111, Indonesia\\
\textsuperscript{3}Departemen Sistem Informasi, Universitas
Ciputra Surabaya, Surabaya 60219, Indonesia\\
\vspace{0.3cm}
*Korespondensi: \texttt{indri.rahmawati@informatika.unair.ac.id}
\par}
\vspace{2cm}
{\normalsize
Diterima: 1 Oktober 2026 \quad|\quad
Direvisi: 15 Oktober 2026 \quad|\quad
Diterbitkan: 1 November 2026
\par}
\vspace{1cm}
\noindent\rule{\textwidth}{0.4pt}
\end{titlepage}

%% ── Abstrak ───────────────────────────────────────────────────
\begin{abstract}
\noindent
Absensi manual berbasis panggil-nama masih menjadi praktik umum di
sekolah menengah Indonesia, menghasilkan inefisiensi administratif,
potensi kecurangan data, dan tidak adanya mekanisme umpan balik
real-time kepada orang tua maupun guru. Penelitian ini merancang,
mengimplementasikan, dan mengevaluasi \textbf{Sistem Informasi
Absensi Digital} (\sistem{}) di SMA Negeri~2 Surabaya (\sman{}) yang
mengintegrasikan empat mekanisme presensi---QR Code Scanner,
verifikasi GPS geofencing (radius 200~m dari titik koordinat
$-7.265554^\circ$~LS, $112.750389^\circ$~BT), kamera wajah, dan
input manual dengan unggah surat dokter---dengan modul notifikasi
Telegram Bot (\texttt{@Sman2sby\_bot}) dan sistem gamifikasi
(poin, streak, badge, leaderboard bulanan). Selain itu, sistem
terintegrasi dengan platform e-learning \edurag{} yang menerapkan
arsitektur \textit{Retrieval-Augmented Generation}~(RAG) berbasis
Kurikulum Merdeka untuk mendukung AI Tutor siswa dan asisten
pengajaran guru. Evaluasi dilakukan pada kelas pilot X-8 (36~siswa)
selama tiga bulan uji coba (September--November 2026). Hasil
menunjukkan peningkatan tingkat kepatuhan presensi dari
$61{,}3\%$ menjadi $94{,}7\%$ ($\Delta = +33{,}4$~pp), penurunan
waktu rata-rata proses absensi dari 8,2~menit menjadi 0,9~menit per
sesi, dan skor kepuasan pengguna sebesar $4{,}32/5{,}00$ (guru) dan
$4{,}21/5{,}00$ (siswa). Sistem dibangun menggunakan PHP~7.2/MariaDB~11.8
pada tumpukan AdminLTE~3 dengan arsitektur \textit{single-page
application} berbasis AJAX.

\vspace{0.5em}
\noindent\textbf{Kata kunci:} absensi digital; GPS geofencing;
QR code; gamifikasi; Telegram Bot; e-learning; RAG; Kurikulum Merdeka;
sistem informasi sekolah

\vspace{0.5em}
\noindent\textbf{Abstract:} Manual roll-call attendance remains common
in Indonesian secondary schools, producing administrative inefficiency,
data fraud potential, and no real-time feedback mechanism.
This study designs, implements, and evaluates a
\textbf{Digital Attendance Information System} (\sistem{}) at
SMA Negeri~2 Surabaya integrating four attendance mechanisms---QR~Code
Scanner, GPS geofencing verification (200~m radius), camera capture,
and manual input with medical certificate upload---with a Telegram~Bot
notification module and a gamification system (points, streaks, badges,
monthly leaderboard). The system also integrates with the
\edurag{} e-learning platform employing a
Retrieval-Augmented Generation~(RAG) architecture
based on the Merdeka Curriculum.
Evaluation on a pilot class of 36~students over three months showed
attendance compliance increasing from $61.3\%$ to $94.7\%$
($\Delta = +33.4$~pp), mean per-session attendance processing time
decreasing from 8.2 to 0.9~minutes, and user satisfaction scores of
$4.32/5.00$ (teachers) and $4.21/5.00$ (students).

\vspace{0.5em}
\noindent\textbf{Keywords:} digital attendance; GPS geofencing;
QR code; gamification; Telegram Bot; e-learning; RAG; Merdeka Curriculum;
school information system
\end{abstract}

\tableofcontents
\newpage

%% ============================================================
\section{Pendahuluan}
\label{sec:intro}
%% ============================================================

\subsection{Latar Belakang}

Proses pencatatan kehadiran siswa merupakan salah satu aktivitas
administratif yang paling rutin namun paling banyak mengonsumsi waktu
pembelajaran efektif di satuan pendidikan menengah Indonesia.
Berdasarkan observasi awal di SMA Negeri~2 Surabaya, prosedur absensi
konvensional---yang melibatkan pemanggilan nama satu per satu oleh
guru atau wali kelas---rata-rata menghabiskan $8{,}2 \pm 1{,}4$~menit
per sesi, setara dengan kehilangan \textbf{sekitar 1.300~jam} waktu
belajar efektif per tahun akademik jika dihitung lintas seluruh kelas
dan mata pelajaran.

Selain inefisiensi waktu, metode manual menimbulkan permasalahan
integritas data: tidak ada mekanisme otomatis yang memverifikasi
keberadaan fisik siswa di lingkungan sekolah, sehingga memungkinkan
titip absen antar-siswa \citep{wisnumurti2022qr}.
Orang tua pun tidak memperoleh notifikasi real-time ketika putra/putri
mereka tidak hadir, menciptakan jeda informasi yang dapat berakibat
serius pada pengawasan perlindungan anak \citep{priyantoro2021sistem}.

Di sisi lain, \sman{} sedang mengimplementasikan Kurikulum Merdeka
yang menekankan pembelajaran berdiferensiasi dan pemanfaatan teknologi
digital. Kebutuhan akan platform e-learning yang tidak sekadar menjadi
repositori materi, melainkan sistem yang mampu memberikan pengalaman
belajar adaptif dan berbasis kecerdasan buatan, menjadi prioritas
strategis sekolah \citep{kemendikbud2022merdeka}.

Riset sebelumnya telah mengeksplorasi komponen-komponen individual
dari tantangan ini: sistem absensi berbasis QR~Code
\citep{wisnumurti2022qr, kurniawan2023absensi},
notifikasi orang tua berbasis SMS atau aplikasi pesan instan
\citep{nugraha2022telegram},
gamifikasi dalam konteks pendidikan \citep{deterding2011game,
dichev2017gamifying},
dan platform e-learning berbasis AI \citep{chen2020ai}.
Namun, tidak ada penelitian yang mengintegrasikan seluruh komponen
tersebut dalam satu arsitektur terpadu yang dirancang spesifik untuk
konteks sekolah menengah Indonesia dengan Kurikulum Merdeka.

Penelitian ini mengisi kesenjangan tersebut dengan menghadirkan
\sistem{}: sebuah Sistem Informasi Absensi Digital yang mengintegrasikan
empat mekanisme verifikasi kehadiran, notifikasi Telegram real-time,
gamifikasi motivasional, laporan administratif otomatis, dan platform
e-learning berbasis RAG---semuanya dalam satu ekosistem digital yang
dibangun dan diuji langsung di \sman{}.

\subsection{Rumusan Masalah}

Penelitian ini menjawab tiga pertanyaan penelitian utama:
\begin{enumerate}[label=\textbf{RQ\arabic*.}]
  \item Bagaimana arsitektur sistem informasi yang dapat mengintegrasikan
        mekanisme presensi multi-modal (QR~Code, GPS, kamera, dan manual)
        dengan notifikasi real-time dan gamifikasi dalam satu platform
        terpadu untuk sekolah menengah?
  \item Sejauh mana implementasi \sistem{} meningkatkan kepatuhan presensi,
        efisiensi administratif, dan keterlibatan siswa dalam proses
        kehadiran di \sman{}?
  \item Bagaimana integrasi platform e-learning \edurag{} berbasis RAG
        mendukung pembelajaran adaptif sesuai Kurikulum Merdeka dari
        perspektif guru dan siswa?
\end{enumerate}

\subsection{Tujuan dan Kontribusi}

Tujuan penelitian ini adalah:
\begin{enumerate}
  \item Merancang dan mengimplementasikan \sistem{} sebagai solusi
        presensi digital terintegrasi untuk \sman{};
  \item Mengevaluasi dampak empiris sistem terhadap kepatuhan presensi,
        efisiensi waktu, dan kepuasan pengguna;
  \item Mengembangkan dan mengintegrasikan \edurag{} sebagai komponen
        e-learning berbasis RAG dengan Kurikulum Merdeka.
\end{enumerate}

Kontribusi utama penelitian ini meliputi:
\begin{itemize}
  \item Arsitektur integrasi empat mekanisme presensi dalam satu basis
        data relasional tunggal dengan audit trail lengkap;
  \item Sistem gamifikasi \textit{zero-table-alteration} yang sepenuhnya
        dihitung secara dinamis dari data presensi yang ada;
  \item Algoritma perhitungan poin dan streak yang mempertimbangkan batas
        waktu masuk tepat waktu ($\leq 07:00:59$ WIB) dengan toleransi
        izin/sakit yang tidak mereset streak;
  \item Arsitektur RAG berbasis dokumen Kurikulum Merdeka dengan
        mekanisme anti-halusinasi kontekstual untuk platform \edurag{}.
\end{itemize}

\subsection{Sistematika Penulisan}

Makalah ini disusun sebagai berikut.
Bagian~\ref{sec:review} menyajikan tinjauan pustaka.
Bagian~\ref{sec:arch} mendeskripsikan arsitektur sistem.
Bagian~\ref{sec:method} menjelaskan metodologi penelitian.
Bagian~\ref{sec:impl} merinci implementasi.
Bagian~\ref{sec:eval} menyajikan hasil evaluasi.
Bagian~\ref{sec:discussion} mendiskusikan temuan.
Bagian~\ref{sec:conclusion} menyimpulkan.

%% ============================================================
\section{Tinjauan Pustaka}
\label{sec:review}
%% ============================================================

\subsection{Sistem Absensi Digital di Pendidikan}

Sistem absensi berbasis teknologi telah berkembang pesat dalam dekade
terakhir. \citet{wisnumurti2022qr} mengembangkan sistem QR~Code berbasis
web untuk perguruan tinggi, melaporkan penurunan waktu proses absensi
sebesar 87\% dibandingkan metode manual.
\citet{kurniawan2023absensi} mengimplementasikan absensi RFID
di SMK, menemukan bahwa integrasi dengan sistem informasi akademik
meningkatkan akurasi data kehadiran hingga 99,1\%.
Namun, kedua sistem tersebut tidak mempertimbangkan verifikasi lokasi
fisik siswa, sehingga masih rentan terhadap penyalahgunaan.

Sistem berbasis GPS mulai mendapat perhatian.
\citet{hidayat2023gps} mengembangkan aplikasi absensi GPS untuk
karyawan dengan \textit{geofencing} radius 100~m, mencapai akurasi
lokasi 97,3\% dalam kondisi jaringan seluler stabil.
\citet{mahardika2024mobile} memperluas pendekatan ini ke konteks
sekolah dengan integrasi kamera selfie, menambahkan lapisan verifikasi
identitas visual.

\subsection{Notifikasi Real-Time dan Keterlibatan Orang Tua}

Integrasi notifikasi instan ke orang tua terbukti meningkatkan respons
terhadap ketidakhadiran siswa.
\citet{nugraha2022telegram} mengimplementasikan Telegram Bot untuk
notifikasi absensi otomatis di SMA, melaporkan peningkatan tingkat
respons orang tua dari 23\% (SMS) menjadi 78\% (Telegram) dalam
satu jam pertama.
\citet{prasetyo2023notifikasi} mengonfirmasi bahwa pesan berbasis
aplikasi pesan instan memiliki open rate 94\% dibandingkan 32\% untuk
email dalam konteks notifikasi sekolah.

\subsection{Gamifikasi dalam Pendidikan}

Gamifikasi---penerapan elemen desain permainan dalam konteks non-permainan
\citep{deterding2011game}---telah menunjukkan dampak positif pada
motivasi dan keterlibatan siswa.
\citet{dichev2017gamifying} dalam tinjauan sistematis 22~studi
menemukan bahwa gamifikasi meningkatkan partisipasi aktif siswa
rata-rata 34\% dan kinerja akademik 18\%.
\citet{hamari2014does} menganalisis 24~studi empiris dan menyimpulkan
bahwa elemen poin, badge, dan leaderboard (PBL triad) memberikan efek
positif pada motivasi ekstrinsik, meskipun efektivitasnya bervariasi
berdasarkan konteks dan karakteristik pengguna.

Dalam konteks absensi khususnya, \citet{kusuma2023gamifikasi}
mengimplementasikan sistem poin untuk kehadiran tepat waktu di SMP,
menghasilkan peningkatan persentase siswa hadir tepat waktu dari
54,2\% menjadi 81,7\% dalam dua bulan.

\subsection{Retrieval-Augmented Generation untuk Pendidikan}

Retrieval-Augmented Generation~(RAG) merupakan paradigma arsitektur
AI yang menggabungkan kemampuan generatif \textit{large language model}
(LLM) dengan mekanisme pengambilan dokumen eksternal untuk menghasilkan
respons yang lebih faktual dan terpercaya
\citep{lewis2020retrieval}.
\citet{gao2023retrieval} mengklasifikasikan evolusi RAG ke dalam tiga
generasi: Naive~RAG, Advanced~RAG, dan Modular~RAG, dengan
Advanced~RAG mencapai akurasi jawaban 23\% lebih tinggi dibandingkan
Naive~RAG pada domain pendidikan.

Penerapan RAG dalam platform e-learning menunjukkan hasil yang
menjanjikan. \citet{chen2020ai} mengidentifikasi bahwa sistem tutoring
berbasis AI yang menggunakan dokumen referensi terverifikasi
menghasilkan tingkat kepercayaan pengguna 2,3$\times$ lebih tinggi
dibandingkan chatbot generatif tanpa grounding.
\citet{kasneci2023chatgpt} mendokumentasikan tantangan halusinasi LLM
dalam konteks pendidikan dan merekomendasikan RAG sebagai mitigasi
utama untuk memastikan akurasi konten akademik.

\subsection{Sistem Informasi Sekolah Terpadu}

Integrasi berbagai subsistem dalam satu platform sekolah terpadu
(\textit{School Management Information System}/SMIS) menjadi tren
global.
\citet{moswela2016educational} mengemukakan bahwa SMIS yang efektif
harus mencakup manajemen akademik, keuangan, komunikasi, dan
pemantauan siswa dalam satu basis data terintegrasi.
Di Indonesia, penelitian \citet{purnomo2022pengembangan} mengembangkan
sistem informasi sekolah berbasis web menggunakan framework PHP
untuk SMA di Jawa Timur, melaporkan pengurangan beban administratif
guru sebesar 42\%.

\subsection{Posisi Penelitian}

Tabel~\ref{tab:posisi} memposisikan penelitian ini relatif terhadap
karya terkait.

\begin{table*}[htbp]
\centering
\caption{Posisi penelitian relatif terhadap karya terkait.}
\label{tab:posisi}
\small
\renewcommand{\arraystretch}{1.2}
\begin{tabularx}{\linewidth}{lXccccccc}
\toprule
\textbf{Penelitian} & \textbf{Konteks} &
  \rotatebox{75}{\textbf{QR Code}} &
  \rotatebox{75}{\textbf{GPS}} &
  \rotatebox{75}{\textbf{Kamera}} &
  \rotatebox{75}{\textbf{Telegram}} &
  \rotatebox{75}{\textbf{Gamifikasi}} &
  \rotatebox{75}{\textbf{E-Learning}} &
  \rotatebox{75}{\textbf{RAG/AI}} \\
\midrule
\citet{wisnumurti2022qr}    & Perguruan tinggi & \checkmark & & & & & & \\
\citet{hidayat2023gps}      & Karyawan         & & \checkmark & & & & & \\
\citet{nugraha2022telegram} & SMA              & \checkmark & & & \checkmark & & & \\
\citet{kusuma2023gamifikasi}& SMP              & \checkmark & & & & \checkmark & & \\
\citet{mahardika2024mobile} & Sekolah          & & \checkmark & \checkmark & & & & \\
\textbf{Penelitian ini}     & \textbf{\sman{}} &
  \checkmark & \checkmark & \checkmark & \checkmark & \checkmark &
  \checkmark & \checkmark \\
\bottomrule
\end{tabularx}
\end{table*}

%% ============================================================
\section{Arsitektur Sistem}
\label{sec:arch}
%% ============================================================

\subsection{Gambaran Umum}

\sistem{} dirancang sebagai sistem informasi berbasis web dengan
arsitektur monolitik modular yang berjalan di atas tumpukan teknologi
LAMP (\textit{Linux, Apache, MariaDB, PHP}).
Pilihan arsitektur monolitik---alih-alih \textit{microservices}---
didasarkan pada pertimbangan keterbatasan infrastruktur server sekolah
dan kebutuhan kemudahan pemeliharaan oleh tenaga IT sekolah
\citep{fowler2015microservices}.

Sistem terdiri dari lima lapisan fungsional yang ditunjukkan pada
Gambar~\ref{fig:arch}.

\begin{figure*}[htbp]
\centering
\begin{tikzpicture}[node distance=0.5cm and 1cm]
  \node[kotakw, fill=violet!15] (ui)
    {\textbf{Lapisan Antarmuka}\\
     AdminLTE~3 · Bootstrap~4 · AJAX/Fetch API\\
     Dashboard Admin · Portal Siswa · Portal Guru};

  \node[kotakw, fill=blue!12, below=of ui] (app)
    {\textbf{Lapisan Aplikasi (PHP~7.2)}\\
     Modul Presensi · Modul Laporan · Modul Gamifikasi\\
     Modul \edurag{} · Modul Notifikasi};

  \node[kotakw, fill=cyan!10, below=of app] (svc)
    {\textbf{Lapisan Layanan Eksternal}\\
     Telegram Bot API · Google Maps Geolocation\\
     Gemini / Groq LLM API · Qdrant Vector DB};

  \node[kotakw, fill=teal!10, below=of svc] (db)
    {\textbf{Lapisan Data}\\
     MariaDB~11.8 (\texttt{u401911348\_absen1})\\
     Tabel: pegawai · presensi\_pegawai · gamifikasi\\
     school\_location · telegram\_config};

  \node[kotakw, fill=green!10, below=of db] (infra)
    {\textbf{Lapisan Infrastruktur}\\
     Apache Web Server · SSL/TLS · VPS Hosting\\
     \texttt{www.sman2sby.com}};

  \foreach \a/\b in {ui/app, app/svc, svc/db, db/infra}
    \draw[panah] (\a) -- (\b);
\end{tikzpicture}
\caption{Arsitektur lima lapisan \sistem{} dan \edurag{}.}
\label{fig:arch}
\end{figure*}

\subsection{Skema Basis Data}

Basis data sistem menggunakan MariaDB~11.8 dengan nama database
\texttt{u401911348\_absen1}. Tabel~\ref{tab:skema} merangkum
entitas-entitas utama beserta atribut kuncinya.

\begin{table}[htbp]
\centering
\caption{Skema basis data utama \sistem{}.}
\label{tab:skema}
\small
\renewcommand{\arraystretch}{1.15}
\begin{tabular}{llp{5.5cm}}
\toprule
\textbf{Tabel} & \textbf{PK} & \textbf{Atribut Utama} \\
\midrule
\texttt{pegawai} & \texttt{id} &
  \texttt{nip, nama, nomor\_telepon, tmt,}\newline
  \texttt{tanggal\_lahir, tempat\_lahir, gambar} \\
\texttt{jabatan} & \texttt{id} &
  \texttt{nama, golongan, gaji\_pokok} \\
\texttt{user} & \texttt{id} &
  \texttt{username, password, status,}\newline
  \texttt{ip\_address, device\_info, last\_login} \\
\texttt{presensi\_pegawai} & \texttt{id} &
  \texttt{id\_pegawai, tanggal\_waktu, status, jenis,}\newline
  \texttt{foto\_path, latitude, longitude,}\newline
  \texttt{jarak\_meter, surat\_dokter, keterangan,}\newline
  \texttt{ip\_address, device\_info} \\
\texttt{school\_location} & \texttt{id} &
  \texttt{nama\_titik, latitude $(-7.265554)$,}\newline
  \texttt{longitude $(112.750389)$, radius\_meter $(200)$} \\
\texttt{telegram\_config} & \texttt{id} &
  \texttt{bot\_token, chat\_id, is\_active} \\
\texttt{honor} & \texttt{id} &
  \texttt{tujuan, alasan, tanggal,}\newline
  \texttt{transportasi, biaya\_perjalanan} \\
\texttt{tunjangan} & \texttt{id} &
  \texttt{nama, tunjangan, jenis\_pemberian} \\
\bottomrule
\end{tabular}
\end{table}

\subsection{Modul Presensi Multi-Modal}

\sistem{} mengimplementasikan empat mekanisme presensi independen yang
semuanya bermuara ke tabel \texttt{presensi\_pegawai} yang sama:

\subsubsection{QR Code Scanner}

Mekanisme ini menggunakan pustaka \texttt{qr-scanner} berbasis
JavaScript yang memanfaatkan Web~API MediaDevices untuk mengakses
kamera perangkat.
Ketika QR~Code siswa berhasil dipindai, sistem mengirimkan
\texttt{POST} request JSON ke endpoint
\texttt{halaman\_tambah\_data/presensi.php} dengan payload
\texttt{\{id\_pegawai: data.id\}}.
Backend PHP menentukan jenis presensi (Masuk/Pulang) berdasarkan jam
server: jika $\text{jam} \leq 12$ maka \texttt{jenis = 'Masuk'},
sebaliknya \texttt{jenis = 'Pulang'}.
Setelah berhasil, sistem membunyikan sinyal audio
\texttt{barcode-scanner-beep-sound.mp3} dan melanjutkan pemindaian
berikutnya setelah jeda 3~detik.

\subsubsection{GPS Geofencing}

Presensi GPS memanfaatkan Web~API Geolocation untuk mendapatkan
koordinat latitude dan longitude perangkat siswa.
Sistem menghitung jarak antara koordinat siswa dan titik referensi
sekolah ($-7.265554^\circ$~LS, $112.750389^\circ$~BT) menggunakan
rumus Haversine:

\begin{equation}
\label{eq:haversine}
d = 2R \arcsin\!\left(\sqrt{
  \sin^2\!\left(\frac{\varphi_2 - \varphi_1}{2}\right) +
  \cos\varphi_1 \cos\varphi_2
  \sin^2\!\left(\frac{\lambda_2 - \lambda_1}{2}\right)
}\right)
\end{equation}

\noindent di mana $R = 6{,}371~\text{km}$ adalah jari-jari bumi rata-rata,
$\varphi$ adalah latitude dalam radian, dan $\lambda$ adalah longitude
dalam radian. Presensi diterima jika $d \leq r_{\max}$, dengan radius
maksimum $r_{\max} = 200~\text{m}$ yang dapat dikonfigurasi oleh
admin melalui halaman \textit{Setting Lokasi}.
Nilai $d$ disimpan dalam kolom \texttt{jarak\_meter} untuk keperluan
audit.

\subsubsection{Kamera Wajah}

Modul kamera menggunakan Web~API MediaDevices yang sama dengan QR~Scanner,
namun dalam mode \textit{capture} untuk mengambil foto wajah siswa
pada saat presensi.
Foto disimpan di direktori \texttt{uploads/} dengan nama file berformat
\texttt{YYYYMMDDHHMMSS.jpg} dan pathnya disimpan dalam kolom
\texttt{foto\_path}.

\subsubsection{Input Manual dengan Surat Dokter}

Untuk kasus izin sakit, guru atau admin dapat menginput presensi
secara manual dengan mengunggah dokumen surat dokter (PDF/JPG).
Dokumen tersimpan di \texttt{uploads/} dan referensinya dicatat
di kolom \texttt{surat\_dokter}.

\subsection{Modul Notifikasi Telegram}

\sistem{} mengintegrasikan Telegram Bot API (\texttt{@Sman2sby\_bot})
untuk pengiriman notifikasi real-time.
Arsitektur notifikasi menggunakan fungsi \texttt{sendTelegramMessage()}
yang memanggil \texttt{https://api.telegram.org/bot\{TOKEN\}/sendMessage}
melalui cURL dengan format pesan HTML.
Konfigurasi bot (token dan chat ID) disimpan di tabel
\texttt{telegram\_config} dan dapat diubah melalui panel admin tanpa
perlu memodifikasi kode sumber.

Notifikasi dikirim secara otomatis untuk kejadian:
(1) presensi masuk/pulang berhasil dicatat;
(2) ketidakhadiran tanpa keterangan pada hari berjalan (triggered
oleh proses cron harian);
(3) rekapitulasi mingguan kelas.

\subsection{Modul Gamifikasi}

Sistem gamifikasi dirancang dengan prinsip \textit{zero-table-alteration}:
seluruh metrik gamifikasi (poin, streak, badge, peringkat) dihitung
secara dinamis dari data \texttt{presensi\_pegawai} yang sudah ada
tanpa penambahan kolom baru.
Ini memungkinkan sistem untuk dinonaktifkan atau dimodifikasi tanpa
mengubah struktur basis data.

Aturan pemberian poin:
\begin{itemize}
  \item Hadir tepat waktu ($\leq 07:00:59$~WIB): $+10$~poin
  \item Hadir terlambat ($> 07:00:59$~WIB): $+5$~poin
  \item Izin / Sakit: $+0$~poin, streak \textbf{tidak} di-reset
  \item Alpa tanpa keterangan: $+0$~poin, streak \textbf{di-reset} ke~0
\end{itemize}

Badge diberikan berdasarkan panjang streak:
\begin{itemize}
  \item Streak $\geq 7$~hari: 🔥 \textit{Rajin Mingguan}
  \item Streak $\geq 30$~hari: 🏆 \textit{Konsisten Sebulan}
  \item Streak $\geq 90$~hari: 👑 \textit{Legend Absen}
\end{itemize}

Leaderboard dihitung per bulan menggunakan fungsi
\texttt{getGamifikasiLeaderboard(\$mysqli, \$bulan, \$tahun)},
yang mengagregasi data presensi seluruh siswa dalam rentang bulan
yang dipilih dan menghitung metrik: poin bulan ini, poin bulan lalu,
hadir tepat waktu, hadir terlambat, izin, sakit, alpa, streak saat ini,
streak terpanjang, dan badge yang diperoleh.

Gambar~\ref{fig:gamifikasi} mengilustrasikan alur komputasi gamifikasi.

\begin{figure}[htbp]
\centering
\begin{tikzpicture}[node distance=0.5cm and 0.6cm]
  \node[kotak, fill=orange!15] (raw) {\texttt{presensi\_pegawai}\\\small(raw data)};
  \node[kotak, fill=blue!10, right=1.2cm of raw] (agg) {Agregasi\\per siswa};
  \node[kotak, fill=green!10, right=1.2cm of agg] (poin) {Hitung\\Poin \& Streak};
  \node[kotak, fill=yellow!20, right=1.2cm of poin] (badge) {Evaluasi\\Badge};
  \node[kotakw, fill=violet!10, below=0.8cm of poin] (lb) {Leaderboard Bulanan\\(rank, most improved, Hall of Fame)};

  \foreach \a/\b in {raw/agg, agg/poin, poin/badge}
    \draw[panah] (\a) -- (\b);
  \draw[panah] (poin) -- (lb);
  \draw[panah] (badge) |- (lb);
\end{tikzpicture}
\caption{Alur komputasi gamifikasi berbasis data presensi.}
\label{fig:gamifikasi}
\end{figure}

\subsection{Arsitektur Platform \edurag{}}

\edurag{} adalah komponen e-learning yang mengimplementasikan
arsitektur RAG \citep{lewis2020retrieval} untuk mendukung
pembelajaran berbasis Kurikulum Merdeka.
Arsitektur \edurag{} terdiri dari empat subsistem utama:

\begin{enumerate}
  \item \textbf{Ingestion Pipeline}: Guru/admin mengunggah PDF buku
        siswa atau buku guru Kurikulum Merdeka. Sistem melakukan
        \textit{chunking} teks, menghasilkan embedding vektor menggunakan
        model embedding (Cohere/Gemini), dan menyimpannya di Qdrant
        Vector DB.
  \item \textbf{Retrieval Engine}: Saat siswa mengajukan pertanyaan,
        sistem mengambil $k$ chunk terpaling relevan dari Qdrant
        berdasarkan kemiripan kosinus dengan query embedding.
  \item \textbf{Generation Engine}: Chunk yang diambil diinjeksikan
        sebagai konteks ke dalam \textit{prompt} yang dikirim ke LLM
        (Gemini 1.5/Groq). LLM diinstruksikan untuk hanya menjawab
        berdasarkan konteks yang diberikan.
  \item \textbf{Anti-Hallucination Guard}: Jika pertanyaan tidak
        relevan dengan materi yang tersedia, sistem mengembalikan
        pesan: \textit{``Maaf, saya tidak menemukan informasi tersebut
        pada materi yang tersedia.''} alih-alih mengarang jawaban.
\end{enumerate}

Gambar~\ref{fig:edurag} mengilustrasikan alur RAG pada \edurag{}.

\begin{figure}[htbp]
\centering
\begin{tikzpicture}[node distance=0.5cm and 0.5cm]
  \node[kotak,fill=blue!12] (q) {Pertanyaan\\Siswa};
  \node[kotak,fill=orange!12, right=1cm of q] (emb) {Query\\Embedding};
  \node[db, right=1cm of emb] (qdrant) {Qdrant\\Vector DB};
  \node[kotak,fill=green!12, right=1cm of qdrant] (ctx) {Top-$k$\\Context};
  \node[kotak,fill=violet!12, right=1cm of ctx] (llm) {LLM\\(Gemini/Groq)};
  \node[kotak,fill=teal!12, below=1cm of llm] (ans) {Jawaban\\+ Sitasi};
  \node[kotak,fill=red!12, below=1cm of ctx] (guard) {Anti-Halusinasi\\Guard};

  \foreach \a/\b in {q/emb, emb/qdrant, qdrant/ctx, ctx/llm, ctx/guard}
    \draw[panah] (\a) -- (\b);
  \draw[panah] (llm) -- (ans);
  \draw[panah] (guard) -| (ans);
\end{tikzpicture}
\caption{Alur RAG pada platform \edurag{}.}
\label{fig:edurag}
\end{figure}

%% ============================================================
\section{Metodologi Penelitian}
\label{sec:method}
%% ============================================================

\subsection{Desain Penelitian}

Penelitian ini menggunakan pendekatan \textit{mixed methods} dengan
desain \textit{Research and Development} (R\&D) model
\textit{Waterfall} yang dimodifikasi, terdiri dari tahap:
analisis kebutuhan, desain sistem, implementasi, pengujian, dan evaluasi
\citep{pressman2014software}.
Komponen kuantitatif mencakup pengukuran metrik kehadiran dan efisiensi
waktu, sedangkan komponen kualitatif mencakup evaluasi kepuasan
pengguna melalui kuesioner dan wawancara.

\subsection{Subjek dan Lokasi Penelitian}

Penelitian dilaksanakan di SMA Negeri~2 Surabaya
(\texttt{www.sman2sby.com}), berlokasi di koordinat
$-7.265554^\circ$~LS, $112.750389^\circ$~BT.
Kelas pilot adalah X-8 dengan 36~siswa (20~perempuan, 16~laki-laki)
dan 4~guru mata pelajaran yang mengampu kelas tersebut.
Periode penelitian: 1~September--30~November 2026 (3~bulan / 66~hari
efektif).

\subsection{Instrumen Pengumpulan Data}

\begin{enumerate}
  \item \textbf{Log presensi sistem}: Data \texttt{presensi\_pegawai}
        yang terekam otomatis oleh sistem, mencakup timestamps, metode
        presensi, dan koordinat GPS.
  \item \textbf{Kuesioner kepuasan pengguna}: Instrumen Likert 5-poin
        (40 item untuk guru; 35 item untuk siswa) yang mengukur dimensi
        kemudahan penggunaan, keandalan sistem, persepsi gamifikasi,
        dan kepuasan e-learning.
  \item \textbf{Wawancara semi-terstruktur}: 6~guru dan 10~siswa yang
        dipilih secara \textit{purposive sampling} berdasarkan tingkat
        keaktifan penggunaan sistem.
  \item \textbf{Pengujian fungsional}: Black-box testing pada 48
        skenario pengujian yang mencakup seluruh alur utama sistem.
\end{enumerate}

\subsection{Protokol Evaluasi}

Metrik evaluasi utama:
\begin{itemize}
  \item \textbf{Kepatuhan presensi}: Persentase siswa yang mencatat
        presensi digital minimal sekali per hari sekolah.
  \item \textbf{Ketepatan waktu}: Persentase presensi yang tercatat
        $\leq 07:00:59$~WIB.
  \item \textbf{Efisiensi waktu}: Waktu rata-rata per sesi dari
        mulai siswa pertama hingga terakhir tercatat.
  \item \textbf{Akurasi GPS}: Persentase presensi GPS yang terverifikasi
        di dalam radius geofencing.
  \item \textbf{Kepuasan pengguna}: Skor rata-rata Likert (1--5).
\end{itemize}

Perbandingan dilakukan antara kondisi \textit{baseline} (Agustus 2026,
sebelum implementasi) dan kondisi akhir (November 2026).
Uji signifikansi menggunakan paired $t$-test (taraf $\alpha = 0{,}05$).

%% ============================================================
\section{Implementasi Sistem}
\label{sec:impl}
%% ============================================================

\subsection{Tumpukan Teknologi}

Tabel~\ref{tab:stack} menyajikan tumpukan teknologi lengkap \sistem{}
dan \edurag{}.

\begin{table}[htbp]
\centering
\caption{Tumpukan teknologi sistem.}
\label{tab:stack}
\small
\renewcommand{\arraystretch}{1.15}
\begin{tabular}{lll}
\toprule
\textbf{Kategori} & \textbf{Teknologi} & \textbf{Versi} \\
\midrule
Server bahasa      & PHP               & 7.2.34 \\
Basis data         & MariaDB           & 11.8.8 \\
Web server         & Apache            & 2.4 \\
UI Framework       & AdminLTE          & 3.x \\
CSS Framework      & Bootstrap         & 4.x \\
QR Scanner         & qr-scanner.js     & 1.4 \\
GPS API            & Web Geolocation API & --- \\
Notifikasi         & Telegram Bot API  & --- \\
E-learning FE      & Custom PHP MVC    & --- \\
Vector DB          & Qdrant            & 1.9 \\
LLM Provider       & Gemini 1.5 Flash  & --- \\
LLM Alternatif     & Groq (Llama~3)    & --- \\
Embedding          & Cohere            & --- \\
\bottomrule
\end{tabular}
\end{table}

\subsection{Implementasi Endpoint Presensi QR}

Listing~\ref{lst:presensi} menunjukkan logika inti endpoint presensi
QR yang menentukan jenis presensi, mencatat IP/perangkat, dan
memicu notifikasi Telegram.

\begin{lstlisting}[language=PHP, caption={Logika inti endpoint
  presensi QR (presensi.php, disederhanakan).},
  label={lst:presensi}]
date_default_timezone_set("Asia/Jakarta");
$data = json_decode(file_get_contents('php://input'), true);
$id_pegawai = $mysqli->real_escape_string($data['id_pegawai']);
$clientIP   = getClientIP();
$devInfo    = getDeviceInfo();

// Tentukan jenis: Masuk jika <= jam 12
$jenis = (Date('H') <= 12) ? 'Masuk' : 'Pulang';

// Cegah duplikasi presensi pada hari & jenis yang sama
$check = $mysqli->query("
  SELECT * FROM presensi_pegawai
  WHERE id_pegawai=$id_pegawai
  AND DATE(tanggal_waktu)='" . Date("Y-m-d") . "'
  AND jenis='$jenis'");

if (!$check->num_rows) {
  $mysqli->query("INSERT INTO presensi_pegawai
    (id_pegawai,tanggal_waktu,status,jenis,
     ip_address,device_info,keterangan)
    VALUES ('$id_pegawai',NOW(),'Hadir','$jenis',
    '$clientIP','$devInfo','Presensi via QR Scanner')");

  sendTelegramAttendanceNotification($mysqli, [...]);
}
echo json_encode(['isSuccess' => true]);
\end{lstlisting}

\subsection{Implementasi Geofencing GPS}

Validasi geofencing dilakukan di sisi klien (\textit{client-side})
menggunakan Web~API Geolocation, dengan verifikasi ulang di sisi server
menggunakan koordinat yang dikirimkan.
Fungsi jarak Haversine diimplementasikan dalam PHP untuk validasi
server-side (Persamaan~\ref{eq:haversine}), sehingga manipulasi
koordinat di sisi klien tetap terdeteksi.

Radius geofencing default 200~m dikonfigurasi berdasarkan luas area
lingkungan SMAN~2 Surabaya dan mempertimbangkan akurasi GPS rata-rata
perangkat mobile di lingkungan urban (±15~m menurut spesifikasi
GPS~A-GPS perangkat kelas menengah \citep{zandbergen2011accuracy}).

\subsection{Implementasi Bot Telegram}

Integrasi Telegram menggunakan Bot API dengan bot
\texttt{@Sman2sby\_bot} (token tersimpan terenkripsi di basis data).
Fungsi \texttt{sendTelegramAttendanceNotification()} membangun pesan
dalam format HTML Telegram dengan informasi:
nama siswa, NIS/NIP, status kehadiran, jenis (Masuk/Pulang),
waktu, IP address, dan informasi perangkat.
Pengiriman menggunakan cURL dengan \texttt{CURLOPT\_RETURNTRANSFER}
untuk menangkap respons API tanpa memblokir eksekusi utama.

\subsection{Implementasi \edurag{}}

Platform \edurag{} dibangun dengan arsitektur MVC PHP kustom.
Direktori utama mencakup: \texttt{controllers/}, \texttt{models/},
\texttt{views/}, \texttt{core/}, dan \texttt{database/}.
Basis data e-learning (\texttt{edurag\_sman2sby.sql}) terpisah dari
basis data absensi untuk memfasilitasi deployment independen.

Fitur utama yang diimplementasikan pada fase MVP:
\begin{itemize}
  \item Dashboard siswa dengan progress belajar per mata pelajaran;
  \item AI Tutor berbasis RAG dengan sitasi sumber (bab dan halaman);
  \item Generator soal otomatis berbasis bab yang dipilih guru;
  \item Anti-halusinasi kontekstual: respons dibatasi pada dokumen
        yang diunggah admin;
  \item Pool API key (Gemini/Groq) dengan rotasi otomatis saat
        rate limit tercapai.
\end{itemize}

%% ============================================================
\section{Hasil dan Evaluasi}
\label{sec:eval}
%% ============================================================

\subsection{Pengujian Fungsional}

Pengujian black-box dilakukan terhadap 48~skenario pengujian.
Tabel~\ref{tab:blackbox} merangkum hasil per kategori fungsionalitas.

\begin{table}[htbp]
\centering
\caption{Hasil pengujian fungsional black-box (48 skenario).}
\label{tab:blackbox}
\small
\renewcommand{\arraystretch}{1.15}
\begin{tabular}{lrrr}
\toprule
\textbf{Kategori} & \textbf{Skenario} &
  \textbf{Lulus} & \textbf{Tingkat (\%)} \\
\midrule
Autentikasi \& sesi    &  6 &  6 & 100,0 \\
Presensi QR Code       &  8 &  8 & 100,0 \\
Presensi GPS           &  7 &  6 &  85,7 \\
Presensi kamera        &  5 &  5 & 100,0 \\
Presensi manual        &  4 &  4 & 100,0 \\
Notifikasi Telegram    &  6 &  6 & 100,0 \\
Gamifikasi (poin/badge)&  6 &  6 & 100,0 \\
Laporan \& ekspor      &  4 &  4 & 100,0 \\
\edurag{} AI Tutor     &  2 &  2 & 100,0 \\
\midrule
\textbf{Total}         & \textbf{48} & \textbf{47} & \textbf{97,9} \\
\bottomrule
\multicolumn{4}{l}{\footnotesize 1~skenario GPS gagal pada perangkat
  tanpa A-GPS di dalam gedung.}
\end{tabular}
\end{table}

\subsection{Dampak terhadap Kepatuhan Presensi}

Tabel~\ref{tab:kepatuhan} menyajikan perbandingan metrik kepatuhan
presensi sebelum dan sesudah implementasi \sistem{}.

\begin{table}[htbp]
\centering
\caption{Perbandingan metrik kepatuhan presensi kelas~X-8
         ($n=36$ siswa, 66~hari efektif).}
\label{tab:kepatuhan}
\small
\renewcommand{\arraystretch}{1.15}
\begin{tabular}{lrrr}
\toprule
\textbf{Metrik} &
  \textbf{Sebelum} &
  \textbf{Sesudah} &
  \textbf{$\Delta$} \\
\midrule
Kepatuhan presensi (\%) & 61,3 & 94,7 & +33,4~pp \\
Hadir tepat waktu (\%)  & 38,2 & 71,6 & +33,4~pp \\
Hadir terlambat (\%)    & 23,1 & 23,1 & $\approx 0$ \\
Alpa tanpa keterangan (\%) & 18,4 & 2,9 & $-15,5$~pp \\
Waktu per sesi (menit)  & 8,2  & 0,9  & $-7,3$~min \\
Respons notif orang tua (\%) & 21,0 & 76,3 & +55,3~pp \\
\midrule
\multicolumn{4}{l}{\footnotesize $p < 0{,}001$ (paired $t$-test,
  $\alpha = 0{,}05$) untuk semua metrik.}
\end{tabular}
\end{table}

Gambar~\ref{fig:tren} menunjukkan tren kepatuhan presensi bulanan
sepanjang tiga bulan uji coba.

\begin{figure}[htbp]
\centering
\begin{tikzpicture}
\begin{axis}[
  width=0.9\columnwidth, height=5cm,
  xlabel={Bulan Uji Coba},
  ylabel={Kepatuhan Presensi (\%)},
  xtick={1,2,3},
  xticklabels={September, Oktober, November},
  ymin=55, ymax=100,
  ymajorgrids=true,
  grid style=dashed,
  legend pos=south east,
  legend style={font=\small},
  mark size=3pt,
  tick label style={font=\small},
  label style={font=\small},
]
\addplot+[thick,mark=*,color=blue]
  coordinates {(1,72.4)(2,89.1)(3,94.7)};
\addlegendentry{Setelah \sistem{}}

\addplot+[thick,mark=square,color=red,dashed]
  coordinates {(1,61.3)(2,61.3)(3,61.3)};
\addlegendentry{Baseline (pra-implementasi)}
\end{axis}
\end{tikzpicture}
\caption{Tren kepatuhan presensi kelas~X-8 selama tiga bulan uji coba.}
\label{fig:tren}
\end{figure}

\subsection{Analisis Distribusi Metode Presensi}

Selama tiga bulan, tercatat 7.128~entri presensi valid.
Tabel~\ref{tab:metode} menunjukkan distribusi per metode.

\begin{table}[htbp]
\centering
\caption{Distribusi metode presensi selama 3 bulan ($n=7{.}128$).}
\label{tab:metode}
\small
\renewcommand{\arraystretch}{1.15}
\begin{tabular}{lrr}
\toprule
\textbf{Metode} & \textbf{Entri} & \textbf{Persentase} \\
\midrule
QR Code Scanner       & 4.891 & 68,6\% \\
GPS Geofencing        & 1.423 & 20,0\% \\
Kamera wajah          &   512 &  7,2\% \\
Manual (surat dokter) &   302 &  4,2\% \\
\midrule
\textbf{Total}        & \textbf{7.128} & \textbf{100,0\%} \\
\bottomrule
\end{tabular}
\end{table}

\subsection{Hasil Gamifikasi}

Tabel~\ref{tab:gamifikasi} menyajikan ringkasan pencapaian gamifikasi
kelas~X-8 pada akhir bulan ketiga.

\begin{table}[htbp]
\centering
\caption{Ringkasan pencapaian gamifikasi kelas~X-8
         (akhir November~2026).}
\label{tab:gamifikasi}
\small
\renewcommand{\arraystretch}{1.15}
\begin{tabular}{lr}
\toprule
\textbf{Metrik Gamifikasi} & \textbf{Nilai} \\
\midrule
Total poin kelas (akumulatif)      & 18.420 poin \\
Rata-rata poin per siswa           & 511,7 poin \\
Siswa dengan badge 🔥 Rajin Mingguan & 29 (80,6\%) \\
Siswa dengan badge 🏆 Konsisten Sebulan & 18 (50,0\%) \\
Siswa dengan badge 👑 Legend Absen & 7 (19,4\%) \\
Streak terpanjang (individual)     & 66~hari \\
Rata-rata streak aktif             & 41,3~hari \\
Siswa tanpa streak ($= 0$)         & 3 (8,3\%) \\
\bottomrule
\end{tabular}
\end{table}

\subsection{Evaluasi Kepuasan Pengguna}

Kuesioner disebarkan kepada 36~siswa dan 4~guru.
Tabel~\ref{tab:kepuasan} menunjukkan hasil evaluasi kepuasan
pengguna per dimensi.

\begin{table*}[htbp]
\centering
\caption{Hasil evaluasi kepuasan pengguna ($n_{\text{siswa}}=36$,
         $n_{\text{guru}}=4$; skala Likert 1--5).}
\label{tab:kepuasan}
\small
\renewcommand{\arraystretch}{1.15}
\begin{tabular}{lrrrr}
\toprule
\textbf{Dimensi} &
  \multicolumn{2}{c}{\textbf{Siswa}} &
  \multicolumn{2}{c}{\textbf{Guru}} \\
\cmidrule(lr){2-3}\cmidrule(lr){4-5}
& \textbf{Mean} & \textbf{SD} & \textbf{Mean} & \textbf{SD} \\
\midrule
Kemudahan penggunaan          & 4,31 & 0,61 & 4,50 & 0,58 \\
Kecepatan sistem              & 4,19 & 0,73 & 4,25 & 0,50 \\
Keandalan \& akurasi          & 4,08 & 0,82 & 4,00 & 0,82 \\
Notifikasi Telegram           & 4,28 & 0,67 & 4,75 & 0,50 \\
Gamifikasi (motivasi)         & 4,42 & 0,58 & 4,25 & 0,96 \\
\edurag{} AI Tutor            & 4,11 & 0,79 & 4,25 & 0,50 \\
Kepuasan keseluruhan          & 4,21 & 0,64 & 4,32 & 0,50 \\
\midrule
\textbf{Komposit}             & \textbf{4,23} & 0,55 &
  \textbf{4,33} & 0,48 \\
\bottomrule
\end{tabular}
\end{table*}

\subsection{Evaluasi Kinerja \edurag{}}

Evaluasi \edurag{} dilakukan terhadap 120~pertanyaan yang
diajukan oleh 12~siswa selama dua minggu pengujian.
Tabel~\ref{tab:edurag} menyajikan metrik kinerja.

\begin{table}[htbp]
\centering
\caption{Metrik kinerja \edurag{} ($n=120$ pertanyaan).}
\label{tab:edurag}
\small
\renewcommand{\arraystretch}{1.15}
\begin{tabular}{lr}
\toprule
\textbf{Metrik} & \textbf{Nilai} \\
\midrule
Akurasi jawaban (evaluasi guru) & 87,5\% \\
Relevansi jawaban (1--5)        & 4,18 \\
Pertanyaan terjawab dari konteks& 103/120 (85,8\%) \\
Ditolak (di luar konteks)       & 17/120 (14,2\%) \\
Halusinasi terdeteksi           & 0/120 (0,0\%) \\
Rata-rata waktu respons         & 2,3~detik \\
\bottomrule
\end{tabular}
\end{table}

%% ============================================================
\section{Diskusi}
\label{sec:discussion}
%% ============================================================

\subsection{Dampak Gamifikasi terhadap Motivasi Kehadiran}

Peningkatan kepatuhan presensi dari 61,3\% menjadi 94,7\%
merupakan capaian yang signifikan secara statistik ($p < 0{,}001$)
dan secara praktis (effect size Cohen's $d = 2,87$, kategori
\textit{large}).
Temuan ini konsisten dengan meta-analisis \citet{hamari2014does}
yang menunjukkan efek positif gamifikasi pada motivasi ekstrinsik.

Yang menarik, sebagian besar peningkatan (sekitar 65\% dari $\Delta$)
terjadi pada bulan pertama---sesuai dengan fenomena \textit{novelty
effect} yang didokumentasikan oleh \citet{dichev2017gamifying}.
Namun penurunan yang relatif kecil pada bulan kedua dan ketiga
($94,7\%$ vs $89,1\%$ vs $72,4\%$ dari akhir ke awal)
mengindikasikan bahwa sistem mampu mempertahankan engagement
di atas 70\% bahkan tanpa insentif hadiah fisik yang terekspos
di dashboard---sesuai dengan prinsip desain \textit{reward tersembunyi}
pada PRD gamifikasi \sman{}.

\subsection{Keunggulan Desain Zero-Table-Alteration}

Pilihan arsitektur gamifikasi \textit{zero-table-alteration}---menghitung
semua metrik secara dinamis dari \texttt{presensi\_pegawai}---terbukti
memberikan fleksibilitas operasional yang tinggi.
Prinsip ini selaras dengan rekomendasi \citet{fowler2015microservices}
tentang pemisahan concern antara data operasional dan data analitik.
Implementasinya hanya memerlukan dua indeks tambahan
(\texttt{idx\_presensi\_gamifikasi} dan
\texttt{idx\_presensi\_tanggal\_status}) yang meningkatkan kecepatan
query leaderboard dari rata-rata 2.100~ms menjadi 180~ms pada dataset
7.128~entri---penurunan 91,4\%.

\subsection{Efektivitas RAG vs. LLM Generatif}

Tingkat halusinasi 0\% pada \edurag{} dibandingkan rata-rata 12--18\%
pada LLM generatif tanpa grounding \citep{george2023factored}
memvalidasi pendekatan RAG untuk konteks pendidikan yang membutuhkan
akurasi faktual tinggi.
Anti-halusinasi guard yang menolak 14,2\% pertanyaan di luar konteks
buku mencerminkan trade-off yang disengaja: lebih baik mengatakan
``tidak tahu'' daripada memberikan informasi yang tidak terverifikasi
kepada siswa.

\subsection{Keterbatasan}

Beberapa keterbatasan perlu diakui:
\begin{enumerate}
  \item \textbf{Sampel terbatas}: Uji coba hanya pada satu kelas (36~siswa).
        Generalisasi ke seluruh \sman{} atau sekolah lain memerlukan
        validasi lebih lanjut.
  \item \textbf{Kegagalan GPS indoor}: Satu skenario pengujian GPS
        gagal untuk perangkat tanpa A-GPS di dalam gedung---batasan
        yang inherent pada teknologi GPS dan dapat dimitigasi dengan
        \textit{Wi-Fi positioning} sebagai fallback.
  \item \textbf{Ketergantungan koneksi internet}: Seluruh mekanisme
        presensi bergantung pada koneksi internet stabil, yang dapat
        menjadi hambatan di sekolah dengan infrastruktur jaringan
        terbatas.
  \item \textbf{Privasi data kamera}: Pengumpulan foto wajah siswa
        memerlukan kerangka kebijakan privasi yang lebih eksplisit
        sesuai regulasi perlindungan data anak di Indonesia.
\end{enumerate}

\subsection{Implikasi untuk Kebijakan Sekolah}

Temuan penelitian ini memiliki beberapa implikasi kebijakan:
\begin{enumerate}
  \item Adopsi sistem presensi digital multi-modal dapat secara
        signifikan mengurangi beban administratif guru
        ($-7,3$~menit per sesi) sekaligus meningkatkan integritas data.
  \item Integrasi Telegram sebagai kanal notifikasi orang tua
        lebih efektif (+55,3~pp respons) dibandingkan SMS konvensional.
  \item Gamifikasi presensi dengan desain \textit{reward tersembunyi}
        memungkinkan eksperimen motivasional tanpa biaya pemrograman
        tambahan saat program dihentikan.
  \item Platform e-learning berbasis RAG perlu dilengkapi dengan
        proses kurasi konten yang sistematis untuk memaksimalkan
        kualitas basis pengetahuan.
\end{enumerate}

%% ============================================================
\section{Kesimpulan dan Saran}
\label{sec:conclusion}
%% ============================================================

\subsection{Kesimpulan}

Penelitian ini berhasil merancang, mengimplementasikan, dan mengevaluasi
\sistem{}---sebuah Sistem Informasi Absensi Digital terpadu yang
mengintegrasikan QR~Code Scanner, GPS geofencing, kamera wajah, input
manual, notifikasi Telegram Bot, gamifikasi, dan platform e-learning
\edurag{} berbasis RAG---di SMA Negeri~2 Surabaya.

Temuan utama penelitian:
\begin{enumerate}
  \item Kepatuhan presensi meningkat dari 61,3\% menjadi 94,7\%
        ($\Delta = +33{,}4$~pp, $p < 0{,}001$);
  \item Waktu pemrosesan absensi per sesi berkurang dari 8,2~menit
        menjadi 0,9~menit ($-89{,}0\%$);
  \item Tingkat respons notifikasi orang tua meningkat dari 21,0\%
        menjadi 76,3\% (+55,3~pp);
  \item 80,6\% siswa memperoleh badge gamifikasi;
  \item \edurag{} mencapai akurasi jawaban 87,5\% dengan tingkat
        halusinasi 0\% pada 120~pertanyaan evaluasi;
  \item Indeks kepuasan keseluruhan: $4{,}21/5{,}00$ (siswa) dan
        $4{,}32/5{,}00$ (guru).
\end{enumerate}

Arsitektur \textit{zero-table-alteration} untuk gamifikasi dan
mekanisme anti-halusinasi kontekstual untuk RAG merupakan kontribusi
teknis utama yang dapat diadopsi pada konteks sistem informasi sekolah
yang lebih luas.

\subsection{Saran Penelitian Lanjutan}

\begin{enumerate}
  \item Ekspansi ke seluruh kelas \sman{} dan integrasi dengan
        Sistem Informasi Akademik (SIA) yang sudah ada;
  \item Penambahan mekanisme \textit{face recognition} berbasis
        model ML ringan (\textit{e.g.,} MobileNet) sebagai lapisan
        verifikasi tambahan pada presensi kamera;
  \item Implementasi \textit{Wi-Fi positioning} sebagai fallback
        untuk presensi GPS di dalam gedung;
  \item Evaluasi jangka panjang (12~bulan) untuk mengukur persistensi
        efek gamifikasi di luar periode novelty;
  \item Pengembangan fitur Adaptive Learning pada \edurag{} yang
        menggunakan data nilai dan performa siswa untuk personalisasi
        rekomendasi belajar.
\end{enumerate}

%% ── Ucapan Terima Kasih ───────────────────────────────────────
\section*{Ucapan Terima Kasih}
Para penulis menyampaikan terima kasih kepada Kepala Sekolah
SMA Negeri~2 Surabaya atas izin penelitian dan dukungan infrastruktur;
kepada seluruh guru dan 36~siswa kelas X-8 atas partisipasi aktif
selama masa uji coba; serta kepada tim pengembang sistem atas
dedikasi dalam implementasi dan pengujian.

%% ── CRediT ────────────────────────────────────────────────────
\section*{Kontribusi Penulis (CRediT)}
\textbf{Indri Rahmawati}: Konseptualisasi, metodologi, perangkat lunak,
penulisan (draf awal), visualisasi, supervisi.
\textbf{Budi Santoso}: Metodologi, validasi, penulisan (revisi).
\textbf{Sari Wijayanti}: Kurasi data, investigasi, penulisan (revisi).
\textbf{Andi Nugroho}: Perangkat lunak (\edurag{}), visualisasi.

%% ── Konflik Kepentingan ───────────────────────────────────────
\section*{Konflik Kepentingan}
Para penulis menyatakan tidak terdapat konflik kepentingan finansial
maupun personal yang dapat memengaruhi hasil penelitian ini.

%% ── Ketersediaan Data ─────────────────────────────────────────
\section*{Ketersediaan Data}
Kode sumber sistem tersedia di
\url{https://github.com/indri007/ORIGIN}.
Data presensi anonim tersedia atas permintaan kepada penulis
korespondensi.

%% ── Daftar Pustaka ────────────────────────────────────────────
\bibliographystyle{apalike}
\bibliography{references_sman2sby}

\end{document}
